\section{Análisis}
\label{sec:analisis}

\subsection{Por qué la CNN supera a la MLP}

La diferencia entre la MLP (F1 0.018) y la CNN de referencia (F1 0.578) no es de
grado sino de tipo: la MLP no logra ajustar ni siquiera el entrenamiento
(F1 de entrenamiento 0.019), mientras que la CNN sí lo ajusta y luego
generaliza parcialmente. La causa está en cómo cada red trata la imagen.

\begin{itemize}
\item \textbf{La MLP no comparte pesos.} Aplana la imagen en 37,632 valores y
conecta cada uno con cada neurona oculta: el 99.7\% de sus 9.66 millones de
parámetros está en esa primera capa. Un pétalo en la esquina superior y el mismo
pétalo en el centro activan pesos distintos, así que la red tendría que aprender
cada patrón por separado en cada posición. Con 122 imágenes por clase en
promedio, no hay datos suficientes para eso.
\item \textbf{La CNN codifica supuestos correctos sobre imágenes.} Sus filtros
de $3\times3$ miran vecindarios locales y se aplican en toda la imagen (pesos
compartidos), las capas sucesivas combinan bordes en texturas y texturas en
partes, y el \emph{pooling} global hace que la predicción dependa de qué hay en
la imagen y no de dónde está. Con 20 veces menos parámetros (480,936) aprovecha
mucho mejor los mismos datos.
\item \textbf{La MLP se refugia en las clases grandes.} Solo predice 11 clases y
su F1 por clase sí depende fuertemente del tamaño de la clase (Spearman 0.51,
frente a 0.12 en la CNN): sin poder distinguir flores, lo mejor que encuentra es
favorecer las clases más frecuentes. Además es inestable: con la tasa de
$10^{-3}$ de la CNN colapsa a una sola clase, por lo que necesitó $10^{-4}$.
\end{itemize}

La comparación no aísla por completo el efecto de las convoluciones: las dos redes
difieren también en número de parámetros y la MLP usa una tasa de aprendizaje
diez veces menor. Aun así, una diferencia de 0.56 puntos de F1 macro con la
misma política de entrenamiento, los mismos datos y el mismo aumento no se
explica por la tasa: con la tasa de la CNN, la MLP es todavía peor.

Los regularizadores no son el problema de la MLP: su curva con aumento y dropout
(gris en la Figura~\ref{fig:curvas-modelos}) es todavía peor que la limpia, y su
brecha es nula. Regularizar más una red que infraajusta no la ayudaría.

\subsection{Por qué la CNN más profunda empeoró}

Agregar dos convoluciones (de 480,936 a 1,218,600 parámetros) bajó el F1 de
0.578 a 0.566. Con una sola semilla, una diferencia de 0.012 es pequeña, pero
coincide con el Avance~4 (0.574 contra 0.550) y las curvas explican por qué la
red más grande no mejora:

\begin{itemize}
\item \textbf{Más capacidad sin más datos.} La red ampliada ajusta el
entrenamiento mucho más rápido: en la época 41 su F1 de entrenamiento es 0.794,
contra 0.638 de la referencia, pero su F1 de validación se estanca en 0.564. La
capacidad extra se usa para memorizar detalles del entrenamiento, sobre todo en
las clases con 18--30 imágenes, no para aprender rasgos que generalicen.
\item \textbf{El sobreajuste aparece antes.} Su pérdida de validación toca su
mínimo en la época 34 (frente a la 54 de la referencia) y después sube, y tiene
seis clases con F1 = 0 en vez de cinco.
\item \textbf{La regularización estaba calibrada para la red pequeña.} Ambas
usan dropout 0.5 y decaimiento de pesos de 0.01. Para una red con 2.5 veces más
parámetros, esa misma regularización es relativamente más débil.
\item \textbf{Más profundidad no agrega información a $112\times112$.} Las dos
capas nuevas operan sobre mapas ya reducidos; a esa resolución, los detalles
finos que distinguen flores parecidas (\emph{carnation} contra \emph{rose}) ya se
perdieron en la entrada, y ninguna capa extra los recupera.
\end{itemize}

En resumen, el cuello de botella de este problema no es la capacidad del modelo
sino la cantidad de datos por clase. Por eso los cambios que atacan la falta de
datos, como el aumento, tienen más efecto que hacer la red más grande.

\subsection{Qué aporta la regularización}

Los tres experimentos de regularización apuntan en la misma dirección: quitar
regularización (Sección~\ref{sec:res-regularizacion}) o reducirla (dropout 0.3)
sube el F1 de entrenamiento o lo deja igual, pero baja el de validación y abre
la brecha. Con 12,753 imágenes y 480,936 parámetros, la CNN tiene capacidad de
sobra para memorizar el entrenamiento; la regularización es la que le permite
generalizar. Su efecto conjunto vale 0.045 puntos de F1 macro.

La regularización no elimina el sobreajuste: la referencia conserva una brecha
de 0.235. Las curvas muestran que la brecha se estabiliza alrededor de la época
60, justo cuando la reducción automática de la tasa lleva el aprendizaje casi a
cero. A partir de ahí la red deja de cambiar, en vez de seguir memorizando. Esto
también explica por qué la referencia toma su mejor checkpoint tan tarde (época
91): son ajustes finos con una tasa de $10^{-5}$.

\subsection{Sensibilidad a los hiperparámetros}

La tasa de aprendizaje es, de lejos, el hiperparámetro más sensible: bajarla a
$3\times10^{-4}$ cuesta 0.118 de F1, más que quitar toda la regularización. La
causa es la interacción con la reducción en meseta: una tasa inicial baja hace
que el calendario la lleve al mínimo antes de que la red aprenda lo suficiente.
Con este calendario, $10^{-3}$ es la elección correcta. El dropout es menos
sensible: entre 0.3 y 0.5 la diferencia es de 0.019, a favor de 0.5.

% PENDIENTE: desbalance (class_weight)

\subsection{Limitaciones}

Cada configuración se entrenó con una sola semilla, por lo que diferencias de
menos de unas centésimas de F1 no deben tomarse como concluyentes. La
validación se usó para elegir el checkpoint de cada corrida, lo que sesga sus
métricas ligeramente hacia arriba; los puntajes de Kaggle
(Sección~\ref{sec:kaggle}) son la única medida sobre datos nunca vistos. La
corrida de la CNN ampliada se interrumpió en la época 50. Todas las redes se
entrenaron desde cero a $112\times112$; no se evaluó transferencia de
aprendizaje desde modelos preentrenados.
